\documentclass[ ../main.tex]{subfiles}
\providecommand{\mainx}{..}
\begin{document}
\section{Random approximate Boolean algebra}
\label{sec:bernoulli_boolean_algebra}

The type $\AT{\Bool}$ introduced in \cref{sec:approx_bitset} models approximation of a single bit.
In this section, we extend the treatment to $n$-bit vectors and show that the
component-wise Boolean operations inherit a tractable error model from the
product-type Kronecker factorization of \cref{thm:product_parsimony}.

\subsection{The Boolean algebra on \texorpdfstring{$\BitSet^n$}{B\^{}n}}

The classical Boolean algebra over $n$-bit vectors is the six-tuple
\begin{equation}
\label{eq:bool_algebra}
	\bigl(\,\BitSet^n,\;\wedge,\;\vee,\;\neg,\;\mathbf{1}^n,\;\mathbf{0}^n\,\bigr)\,,
\end{equation}
where $\wedge$, $\vee$, and $\neg$ denote component-wise AND, OR, and NOT,
and $\mathbf{1}^n$ and $\mathbf{0}^n$ are the all-ones and all-zeros vectors.
The type $\BitSet^n$ is isomorphic to the $n$-fold product $\Bool \times \cdots \times \Bool$,
so any type with cardinality $2^n$ admits a Boolean algebra via this identification.
For instance, $\BitSet^n$ can represent the powerset $2^{\Set{U}}$ of any universe
$\Set{U}$ with $|\Set{U}| = n$, giving a complete implementation of set-theoretic
operations (union, intersection, complement) on $\Set{U}$.

\subsection{Approximate bit vectors}

An approximate $n$-bit vector $\tilde{\mathbf{x}} \in \AT{\BitSet^n}$ replaces each
component $x_i \in \BitSet$ by an independent Bernoulli Boolean
$\tilde{x}_i \in \AT{\Bool}$ with per-component error rates
$(\fprate_i, \fnrate_i)$ for $i = 1, \ldots, n$.
Since $\BitSet^n$ is an $n$-fold product type, the Kronecker factorization
(\cref{thm:product_parsimony}) applies directly:

\begin{theorem}[Kronecker factorization of $\AT{\BitSet^n}$]
\label{thm:bitvec_kronecker}
The joint confusion matrix of $\AT{\BitSet^n}$ factors as
\begin{equation}
\label{eq:bitvec_kronecker}
	M_{\BitSet^n}
	= M_1 \otimes M_2 \otimes \cdots \otimes M_n\,,
\end{equation}
where each $M_i$ is the $2 \times 2$ per-component channel matrix
\begin{equation}
	M_i
	= \begin{pmatrix}
		1 - \fprate_i & \fprate_i \\
		\fnrate_i     & 1 - \fnrate_i
	  \end{pmatrix}\,.
\end{equation}
A general channel on $2^n$ symbols has $2^n(2^n - 1)$ free parameters.
Under the Bernoulli independence assumption, the factorization reduces this to $2n$ parameters:
two rates $(\fprate_i, \fnrate_i)$ per component.
\end{theorem}
\begin{proof}
The type $\BitSet^n = \Bool^n$ is the $n$-fold product $\Bool \times \cdots \times \Bool$.
By \cref{thm:product_parsimony}, the joint confusion matrix factors as the
Kronecker product of the per-component matrices, and the total parameter count
is $\sum_{i=1}^{n} \mathrm{params}(\Bool) = 2n$.
\end{proof}

\begin{corollary}[Correctness of approximate bit vectors]
\label{cor:bitvec_correct}
The probability that an approximate $n$-bit vector is fully correct is
\begin{equation}
	\Prob{\tilde{\mathbf{x}} = \mathbf{x}}
	= \prod_{i=1}^{n}
	  \bigl(x_i\,(1 - \fnrate_i) + (1 - x_i)\,(1 - \fprate_i)\bigr)\,.
\end{equation}
\end{corollary}

\subsection{Approximate Boolean operations}

The Boolean operations $\wedge$, $\vee$, and $\neg$ act component-wise
on $\BitSet^n$, so their approximate counterparts inherit errors from the
component channels.

\begin{theorem}[Approximate AND on Bernoulli Booleans]
\label{thm:approx_and}
Let $\tilde{a}, \tilde{b} \in \AT{\Bool}$ be independent Bernoulli Booleans
with rates $(\fprate_1, \fnrate_1)$ and $(\fprate_2, \fnrate_2)$ respectively.
The output $\tilde{c} = \tilde{a} \wedge \tilde{b}$ has the conditional
error rates
\begin{align}
\label{eq:and_rates}
	\Prob{\tilde{c} = 1 \Given a = 0, b = 0}
		&= \fprate_1 \, \fprate_2\,, \\
	\Prob{\tilde{c} = 1 \Given a = 0, b = 1}
		&= \fprate_1 \, (1 - \fnrate_2)\,, \\
	\Prob{\tilde{c} = 1 \Given a = 1, b = 0}
		&= (1 - \fnrate_1) \, \fprate_2\,, \\
	\Prob{\tilde{c} = 1 \Given a = 1, b = 1}
		&= (1 - \fnrate_1) \, (1 - \fnrate_2)\,.
\end{align}
The joint model on the input pair $(a, b) \in \BitSet^2$ has four partition
blocks (one per input combination), yielding four distinct conditional error
probabilities for the Boolean output, but independence yields only
$4$ free parameters by the Kronecker factorization of $M_1 \otimes M_2$.
\end{theorem}
\begin{proof}
By independence of $\tilde{a}$ and $\tilde{b}$, the joint probability
factors:
$\Prob{\tilde{a} = \alpha, \tilde{b} = \beta \Given a, b}
= \Prob{\tilde{a} = \alpha \Given a} \cdot \Prob{\tilde{b} = \beta \Given b}$.
Since AND is deterministic, $\tilde{c} = 1$ if and only if
$\tilde{a} = 1$ and $\tilde{b} = 1$.
The four cases follow by substituting the per-component channel probabilities.
\end{proof}

\begin{theorem}[Approximate OR on Bernoulli Booleans]
\label{thm:approx_or}
Under the same setup as \cref{thm:approx_and}, the output
$\tilde{c} = \tilde{a} \vee \tilde{b}$ has the conditional error rates
\begin{align}
	\Prob{\tilde{c} = 1 \Given a = 0, b = 0}
		&= 1 - (1 - \fprate_1)(1 - \fprate_2)\,, \\
	\Prob{\tilde{c} = 1 \Given a = 0, b = 1}
		&= 1 - (1 - \fprate_1)\,\fnrate_2\,, \\
	\Prob{\tilde{c} = 1 \Given a = 1, b = 0}
		&= 1 - \fnrate_1\,(1 - \fprate_2)\,, \\
	\Prob{\tilde{c} = 1 \Given a = 1, b = 1}
		&= 1 - \fnrate_1 \, \fnrate_2\,.
\end{align}
\end{theorem}
\begin{proof}
The event $\tilde{c} = 0$ requires $\tilde{a} = 0$ and $\tilde{b} = 0$,
which by independence has probability
$\Prob{\tilde{a} = 0 \Given a} \cdot \Prob{\tilde{b} = 0 \Given b}$.
The result follows from $\Prob{\tilde{c} = 1} = 1 - \Prob{\tilde{c} = 0}$.
\end{proof}

\begin{remark}[Approximate NOT]
\label{rem:approx_not}
For $\tilde{a} \in \AT{\Bool}[\fprate][\fnrate]$, the negation
$\tilde{c} = \neg\,\tilde{a}$ swaps the roles of false positives and
false negatives: $\tilde{c} \in \AT{\Bool}[\fnrate][\fprate]$.
This is the component-wise analogue of the complement operation on
Bernoulli sets, where complementing swaps $\fprate$ and $\fnrate$.
\end{remark}

\subsection{Component-wise operations on \texorpdfstring{$\AT{\BitSet^n}$}{AT(B\^{}n)}}

Extending the above to $n$-bit vectors is immediate: each Boolean operation
acts independently on each component, so the $i$-th output component
depends only on the $i$-th input components and their respective error rates.

\begin{corollary}[Component-wise AND on $\AT{\BitSet^n}$]
\label{cor:bitvec_and}
For independent approximate bit vectors
$\tilde{\mathbf{a}}, \tilde{\mathbf{b}} \in \AT{\BitSet^n}$,
the component-wise AND $\tilde{\mathbf{c}} = \tilde{\mathbf{a}} \wedge \tilde{\mathbf{b}}$
has the $i$-th component error rates given by \cref{thm:approx_and}
applied to $(\fprate_{a,i}, \fnrate_{a,i})$ and $(\fprate_{b,i}, \fnrate_{b,i})$.
The joint confusion matrix of $\tilde{\mathbf{c}}$ factors as a Kronecker
product of $n$ output-component channels, each determined by the two
input-component channels at position $i$.
\end{corollary}

\begin{remark}[Connection to Bernoulli set operations]
When $\BitSet^n$ represents the powerset of an $n$-element universe $\Set{U}$,
the component-wise AND, OR, and NOT correspond to set intersection, union,
and complement.
The error rates derived above are consistent with the set-theoretic
composition rules for Bernoulli sets~\cite{bernoulliComposition}:
approximate intersection has per-element false positive rate
$\fprate_1 \fprate_2$ and approximate union has per-element false negative
rate $\fnrate_1 \fnrate_2$.
\end{remark}

\end{document}
